\subsection{Opposite influence-AA classifications at identical
spectra}\label{quant:sec-3}

Take the complete graph with \(J_{ij}=1/n\). Then

\[
\mathcal H_Z=-\frac1{2n}\left[\left(\sum_iZ_i\right)^2-nI\right].
\]

Its measured input law is precisely the normalized Curie--Weiss law,
since the additive constant has no effect. At each fixed \(\beta>1\),
the globally bounded degree-one polynomial \(f_n(s)=n^{-1}\sum_i s_i\)
satisfies

\[
I_i^\mu(f_n)=n^{-2},\qquad
\operatorname{Var}_\mu(f_n)\longrightarrow m_\beta^2>0,
\qquad m_\beta=\tanh(\beta m_\beta),\quad m_\beta>0.
\]

The variance assertion is proved in
Section~\ref{proof:curie_weiss}. Hence distributional influence AA
fails. For \(\mathcal H_X\) the measured law is uniform at every
temperature, so uniform AA applies.

Therefore \textbf{local interaction eigenvalues, their ordering relative
to the singlet, and even the complete many-body spectrum together with
the graph do not determine fixed-basis Gibbs-input influence AA}.

The obstruction lies inside the single symmetric strictly two-local
interaction class. It uses the same graph, identical spectra and
interaction norms, commuting Hamiltonians, and no added fields.

\subsubsection{A continuous isospectral family: measured thermal
inputs}\label{quant:heading-5}

The example extends to every unit vector \(u\in\mathbb R^3\):

\[
K_u=-(u\cdot\sigma)\otimes(u\cdot\sigma),\qquad
\mathcal H_u(G)=-\sum_{\{i,j\}}J_{ij}(u\cdot\sigma)_i(u\cdot\sigma)_j.
\]

All these Hamiltonians are related by a common single-qubit rotation and
have identical spectra on the same graph. Let \(c=u_z\), the overlap of
the interaction axis with the fixed measurement axis. If \(S\) is the
classical Ising input, the measured thermal law of \(\mathcal H_u\) is
generated by conditionally independent channels

\[
\Pr[X_i=x\mid S_i=s]=(1+cxs)/2.
\]

This follows by diagonalizing each spin along \(u\): the thermal state
is a mixture over the Ising spin configurations in that product basis,
and the squared measurement overlaps are exactly the displayed
probabilities. For normalized magnetizations,

\[
\mathbb E[M_X\mid S]=cM_S,\qquad
\operatorname{Var}(M_X\mid S)=(1-c^2)/n,
\]

so the law of total variance gives the exact identity

\[
\boxed{\operatorname{Var}(M_X)=c^2\operatorname{Var}(M_S)+(1-c^2)/n.}
\]

At fixed \(\beta>1\) on normalized complete graphs, every fixed
\(c\ne0\) therefore violates influence AA through the degree-one
magnetization, whereas \(c=0\) gives the uniform input law exactly. This
is a complete orientation classification for influence AA in this
ordered, isospectral rank-one interaction subfamily.

The positive result includes the critical point, with constants uniform
over orientations:

\[
\boxed{\max_i I_i^{\mu_{\beta,u}}(f)
\ge c_*\operatorname{Var}_{\mu_{\beta,u}}(f)^4/d^{12}
\quad(0\le\beta\le1,\ |u|=1).}
\]

Here \(c_*>0\) is absolute. To verify this, use the Curie--Weiss latent
representation proved in Section~\ref{proof:curie_weiss}. Let
\(Z\) have density proportional to
\(\exp[-n(z^2/(2\beta)-\log\cosh z)]\), and put
\(b=[n(\beta^{-1}-1)+\sqrt n]^{-1/2}\). The latent law satisfies a
Poincaré inequality with constant at most \(Cb^2\), and \(P=\tanh Z\) is
symmetric, has \(\mathbb EP^2\le Cn^{-1/2}\), and has density at least
\(c_0/b\) on \([-b/2,b/2]\).

Passing \(S\) through the measurement channel makes the measured law
exactly the product mixture with mean \(T=c\tanh Z\). If \(c=0\), it is
uniform. If \(c\ne0\), \(T\) is symmetric,
\(\mathbb ET^2\le Cn^{-1/2}\), and its density is at least
\(c_0/(|c|b)\) on \([-|c|b/2,|c|b/2]\). Thus both the low-degree moment
comparison and central-density influence comparison of Lemmas HT3--HT4
in Section~\ref{proof:high_temperature} hold with the same
absolute constants, regardless of how small \(|c|\) is.

For the spin Poincaré estimate put \(g(z)=\mathbb E_{P_{c\tanh z}}f\).
Conditional product Poincaré and differentiation give

\[
\mathbb E_Z\operatorname{Var}_{P_{c\tanh Z}}f\le\sum_i I_i^{\mu_{\beta,u}}(f),
\] \[
g'(z)=c\,\operatorname{sech}^2z\sum_i\mathbb E_{P_{c\tanh z}}D_if,
\qquad
|g'(z)|^2\le n\sum_i\mathbb E_{P_{c\tanh z}}(D_if)^2.
\]

The latent Poincaré inequality therefore yields
\(\operatorname{Var}_{\mu_{\beta,u}}f\le C\sqrt n\sum_i I_i^{\mu_{\beta,u}}(f)\).
All three inputs to Theorem HT5 hold uniformly in \(c\), so its two-case
proof gives the displayed bound unchanged. At \(\beta=0\) the law is
uniform.

The same classification holds for heat-bath influence. Conditioning on
the underlying Ising spins shows \(|\mathbb E[X_i\mid X_{-i}]|
=|c\,\mathbb E[S_i\mid X_{-i}]|\le |c|\tanh\beta\): given the other
Ising spins, the conditional mean of \(S_i\) has absolute value at most
\(\tanh\beta\), and observing their noisy outputs cannot increase this
uniform bound. Hence \(\operatorname{sech}^2\beta\,I_i\le B_i\le I_i\).

This completes the fixed-parameter diagram for normalized complete
graphs: influence AA holds \textbf{if and only if} \(\beta\le1\) or
\(u_z=0\), for fixed \(\beta\ge0\) and fixed orientation \(u\). It does
not classify arbitrary dimension-dependent approaches to this boundary.

Every measured law in this complete-graph family is exchangeable, so all
orientations and all temperatures satisfy acceptance simulation by the
theorem used in Section~\ref{quant:sec-5}.

\subsubsection{The diagonal-energy model has a different
diagram}\label{quant:heading-6}

For the same \(\mathcal H_u\), taking diagonal energies before
exponentiation gives

\[
\langle x|\mathcal H_u(G)|x\rangle
=-c^2\sum_{\{i,j\}}J_{ij}s_is_j.
\]

On the normalized complete graph its classical law is therefore
Curie--Weiss at effective inverse temperature \(\beta c^2\). Its exact
influence-AA criterion is \textbf{\(\beta c^2\le1\)}, including
equality, by the already proved Curie--Weiss theorem. For measured
thermal inputs the criterion was instead \(\beta\le1\) or \(c=0\). For
example, at \(\beta=2,c=1/2\), the measured model fails influence AA
while the diagonal-energy model satisfies it. Both models are
exchangeable here and hence admit acceptance simulation at every
temperature.

The original \(-ZZ\) versus \(-XX\) counterexample is valid in both
models; the intermediate orientations show why they must nevertheless be
defined separately.

\subsubsection{Complete diagram with signed interaction
strength}\label{quant:heading-7}

Allow the fixed interaction

\[
K_{\lambda,u}=-\lambda(u\cdot\sigma)\otimes(u\cdot\sigma),
\qquad |u|=1,\quad \lambda\in\mathbb R,
\]

on the complete graph with edge weights \(1/n\). Put
\(\gamma=\beta\lambda\) and \(c=u_z\). At each fixed
\(\gamma\in\mathbb R\) and fixed orientation,

\[
\boxed{\begin{aligned}
\text{Measured input:}\quad&
\text{AA holds iff }\gamma\le1\text{ or }c=0,\\
\text{Diagonal-energy input:}\quad&
\text{AA holds iff }\gamma c^2\le1.
\end{aligned}}
\]

Every boundary in these criteria is included in the positive region.
Both flip and heat-bath influence have these classifications. Both input
models admit polynomial classical acceptance simulation throughout the
entire diagram.

The preceding arguments proved the claims for \(\gamma\ge0\). For
\(\gamma=-b<0\), the underlying Ising density is
\(r=e^{-bS^2/(2n)}/Z_{n,-b}\). Jensen gives \(r\le e^{b/2}\). For
\(s=1+1/(2b)\),

\[
\mathbb E_{\mu}r^{-s}
=Z_{n,-b}^{s-1}\mathbb E_Ue^{(s-1)bS^2/(2n)}
\le\sqrt2.
\]

The two-sided moment transfer theorem P5 in
Appendix~\ref{proof:positive_structure} therefore
proves full influence AA at every fixed \(b>0\). Its hypotheses are
preserved by every uniform-preserving channel: if \(r_{\rm out}\) is the
conditional average of \(r\), convexity of \(t^p\) and \(t^{1-s}\)
contracts the forward and inverse density moments. Applying the product
measurement channel proves the measured-input statement uniformly over
orientation. The diagonal-energy law is simply signed Curie--Weiss at
coupling \(\gamma c^2\le0\), so P5 applies there directly.

For heat-bath influence, the conditional mean bound through the
underlying Ising spins gives
\(|\mathbb E[X_i\mid X_{-i}]|\le |c|\tanh|\gamma|\). Thus the
conditional variance factor is bounded below at every fixed coupling,
and the same classifications follow.


\subsubsection{Exact classification with the paper's worst-case graph
quantifier}\label{quant:heading-8}

Now fix any \(\beta>0\) and allow all interaction graphs and positive
polynomially bounded edge weights, as in the Marwaha--Sud problem
family. For the single fixed term \(K_u=-(u\cdot\sigma)^{\otimes2}\), a
uniform influence-AA guarantee over this whole family holds \textbf{if
and only if} \(u_z=0\).

If \(u_z=0\), the measurement channel outputs independent fair bits for
every underlying Ising configuration, on every graph and at every
strength, so uniform AA applies. If \(u_z\ne0\), choose complete-graph
weights \(2/(\beta n)\). The underlying Ising law then has effective
Curie--Weiss inverse temperature two. Its measured magnetization has
variance tending to \(u_z^2m_2^2>0\) and all influences \(n^{-2}\),
proving failure of the universal guarantee. These weights are positive
and polynomially bounded. At \(\beta=0\) every thermal state is
maximally mixed and all orientations satisfy AA.

This is a genuine worst-case interaction classification for an
isospectral subclass, rather than merely contrasting easy and hard graph
instances. It concerns influence AA. No worst-case acceptance-simulation
classification over arbitrary graphs is asserted.

The worst-case orientation criterion also holds for the diagonal-energy
model. For \(c\ne0\), use weights \(2/(\beta c^2n)\), producing
effective Curie--Weiss inverse temperature two; for \(c=0\) all diagonal
energies are zero. The different weights are necessary because the two
ways to form classical inputs are not interchangeable.

The signed interaction admits the same complete worst-case
classification. For fixed \(\beta\ge0,\lambda\in\mathbb R,u\), a uniform
influence-AA guarantee over all positive polynomially bounded weighted
graphs holds \textbf{if and only if \(\beta\lambda=0\) or \(u_z=0\)}, in
either classical-input model.

Only the case \(\lambda<0\), \(\beta>0\), \(c=u_z\ne0\) needs an
additional argument. Use the bipartite 16-regular expander family from
Section~\ref{proof:local_interactions},
putting weight \(6/(16\beta|\lambda|)\) on every edge. Its underlying
Ising law is the antiferromagnetic normalized expander law at effective
inverse temperature six. A bipartite spin gauge converts it to the
proved ferromagnetic counterexample. If \(\sigma_i\) is \(+1\) on one
part and \(-1\) on the other, its staggered magnetization
\(M_{\rm stag}(S)=n^{-1}\sum_i\sigma_iS_i\) has variance bounded away
from zero.

For the measured law, the independent measurement channel gives

\[
\operatorname{Var}\!\left(n^{-1}\sum_i\sigma_iX_i\right)
=c^2\operatorname{Var}(M_{\rm stag}(S))+(1-c^2)/n.
\]

This is a constant-variance degree-one witness with every influence
\(n^{-2}\). For the diagonal-energy model use edge weight
\(6/(16\beta|\lambda|c^2)\) instead, so its classical effective coupling
is the same antiferromagnetic expander coupling. Positive \(\lambda\)
follows by rescaling the complete-graph argument; \(\beta\lambda=0\) or
\(c=0\) gives uniform input. Both influence conventions have the same
classification.

The universal graph quantifier matters: antiferromagnetic normalized
complete graphs satisfy AA at every fixed strength, while the
arbitrary-graph class of that same interaction contains these
counterexamples.
